\subsection{Equivalence}

\newcommand{\norlequiv}[2]{\norwedge{\norimplies{#1}{#2}}{\norimplies{#2}{#1}}} % NOR hacking
\begin{notation}
    \[P \lequiv Q \coloneqq (P \implies Q) \wedge (Q \implies P)\]
    \[[\norlequiv{P}{Q}]\]
\end{notation}

% P=>Q, Q=>P > P<=>Q
\begin{align*}
    \premise{P \implies Q}\\
    \premise{Q \implies P}\\
    \byrule{(P \implies Q) \wedge (Q \implies P)}{\ref{rule:conjunction-introduction}}\\
    \bydefinition{P \lequiv Q}
\end{align*}
\begin{equation}
    \inferrule{P \implies Q \\ Q \implies P}{P \lequiv Q}
    \tag{{$\lequiv$}I}
    \label{rule:equivalence-introduction}
\end{equation}

% P<=>Q > P=>Q, Q=>P
\begin{align*}
    \premise{P \lequiv Q}\\
    \bydefinition{(P \implies Q) \wedge (Q \implies P)}\\
    \byrule{P \implies Q}{\ref{rule:conjunction-elimination}}
\end{align*}
\begin{align*}
    \premise{P \lequiv Q}\\
    \bydefinition{(P \implies Q) \wedge (Q \implies P)}\\
    \byrule{Q \implies P}{\ref{rule:conjunction-elimination}}
\end{align*}
\begin{equation}
    \inferrule{P \lequiv Q}{P \implies Q} \quad \inferrule{P \lequiv Q}{Q \implies P}
    \label{rule:equivalence-unpacking}
\end{equation}

% P/Q, P<=>Q > Q/P
\begin{align*}
    \premise{P}\\
    \premise{P \lequiv Q}\\
    \byrule{P \implies Q}{\ref{rule:equivalence-unpacking}}\\
    \byrule{Q}{\ref{rule:conditional-elimination}}\\
\end{align*}
\begin{align*}
    \premise{Q}\\
    \premise{P \lequiv Q}\\
    \byrule{Q \implies P}{\ref{rule:equivalence-unpacking}}\\
    \byrule{P}{\ref{rule:conditional-elimination}}\\
\end{align*}
\begin{equation}
    \inferrule{P \\ P \lequiv Q}{Q} \quad \inferrule{Q \\ P \lequiv Q}{P}
    \tag{{$\lequiv$}E}
    \label{rule:equivalence-elimination}
\end{equation}

% > P<=>P
\begin{align*}
    \byrule{P \implies P}{\ref{rule:conditional-reflexivity}}\\
    \byrule{P \lequiv P}{\ref{rule:equivalence-introduction}}
\end{align*}
\begin{equation}
    \vdash P \lequiv P
    \label{rule:equivalence-reflexivity}
\end{equation}

% P<=>Q > Q<=>P
\begin{align*}
    \premise{P \lequiv Q}\\
    \byrule{P \implies Q}{\ref{rule:equivalence-elimination}}\\
    \byrule{Q \implies P}{\ref{rule:equivalence-elimination}}\\
    \byrule{Q \lequiv P}{\ref{rule:equivalence-introduction}}
\end{align*}
\begin{equation}
    \inferrule{P \lequiv Q}{Q \lequiv P}
    \label{rule:equivalence-symmetry}
\end{equation}

% P<=>Q, Q<=>R > P<=>R
\begin{align*}
    \premise{P \lequiv Q}\\
    \premise{Q \lequiv R}\\
    \assumption[1]{P}\\
    \byrule[1]{Q}{\ref{rule:equivalence-elimination}}\\
    \byrule[1]{R}{\ref{rule:equivalence-elimination}}\\
    \byrule{P \implies R}{\ref{rule:conditional-introduction}}\\
    \byrule{R \lequiv Q}{\ref{rule:equivalence-symmetry}}\\
    \byrule{Q \lequiv P}{\ref{rule:equivalence-symmetry}}\\
    \assumption[1]{R}\\
    \byrule[1]{Q}{\ref{rule:equivalence-elimination}}\\
    \byrule[1]{P}{\ref{rule:equivalence-elimination}}\\
    \byrule{R \implies P}{\ref{rule:conditional-introduction}}\\
    \byrule{P \lequiv R}{\ref{rule:equivalence-introduction}}
\end{align*}
\begin{equation}
    \inferrule{P \lequiv Q \\ Q \lequiv R}{P \lequiv R}
    \label{rule:equivalence-transitivity}
\end{equation}

% P<=>Q, R<=>S > P|R<=>Q|S
\begin{align*}
    \premise{P \lequiv Q}\\
    \premise{R \lequiv S}\\
    \assumption[1]{P \nor R}\\
    \assumption[2]{Q}\\
    \byrule[2]{P}{\ref{rule:equivalence-elimination}}\\
    \byrule[2]{\bot}{\ref{rule:nor-elimination}}\\
    \assumption[2]{S}\\
    \byrule[2]{R}{\ref{rule:equivalence-elimination}}\\
    \byrule[2]{\bot}{\ref{rule:nor-elimination}}\\
    \byrule[1]{Q \nor S}{\ref{rule:nor-introduction}}\\
    \byrule{(P \nor R) \implies (Q \nor S)}{\ref{rule:conditional-introduction}}\\
    \assumption[1]{Q \nor S}\\
    \assumption[2]{P}\\
    \byrule[2]{Q}{\ref{rule:equivalence-elimination}}\\
    \byrule[2]{\bot}{\ref{rule:nor-elimination}}\\
    \assumption[2]{R}\\
    \byrule[2]{S}{\ref{rule:equivalence-elimination}}\\
    \byrule[2]{\bot}{\ref{rule:nor-elimination}}\\
    \byrule[1]{P \nor R}{\ref{rule:nor-introduction}}\\
    \byrule{(Q \nor S) \implies (P \nor R)}{\ref{rule:conditional-introduction}}\\
    \byrule{(P \nor R) \lequiv (Q \nor S)}{\ref{rule:equivalence-introduction}}
\end{align*}
\begin{equation}
    \inferrule{P \lequiv Q \\ R \lequiv S}{(P \nor R) \lequiv (Q \nor S)}
    \label{rule:equivalence-congruence}
\end{equation}

